Serum biomarkers studied in relation to external tissue support (perivascular/adventitial support, external elastic lamina, or periaortic tissue) in aortic disease between 2000 and 2026.
Serum/plasma biomarkers associated with aortic disease (AAA, TAA, AD) and external/perivascular tissue support (adventitia, PVAT, periaortic tissue), drawn from literature ~2000–2026
Direct studies explicitly linking circulating biomarkers to “external tissue support” (as used in continuum FSI models) are sparse. Most data concern ECM remodeling, inflammation, and proteolysis that affect the adventitia and perivascular adipose tissue (PVAT). The list below groups the most consistently reported markers; many were validated or re-evaluated in proteomic/multi-omic studies through 2025–2026.
Core ECM remodeling & proteolysis markers
MMP-9 (matrix metalloproteinase-9) – elevated; correlates with diameter, growth, rupture risk, and PET positivity; reflects collagen/elastin degradation that weakens media–adventitia interface.
MMP-2, MMP-8, MMP-12, MMP-13 – variably elevated; MMP-12 emerged as a strong predictor of future AA/AD risk in large proteomic panels (detectable years before diagnosis).
TIMP-1 / TIMP-2 (tissue inhibitors of metalloproteinases) – often altered (ratios or absolute levels change); imbalance with MMPs drives wall weakening.
Osteopontin (OPN) and Osteoprotegerin (OPG) – elevated with inflammation, calcification, and rapid growth; linked to adventitial remodeling.
Aggrecan / COL4A1 and other ECM fragments – appear in multi-protein panels associated with presence and growth.
Inflammatory & cytokine markers (frequently linked to PVAT/adventitial inflammation)
IL-6 – consistently elevated; associated with presence, growth, organ dysfunction, and mortality; key PVAT-derived adipokine.
CRP / hs-CRP – nonspecific but prognostic for expansion, rupture risk, and post-intervention outcomes.
IL-4, MCP-1 (CCL2), Oncostatin-M, Neurturin – elevated in circulating extracellular vesicles from AAA patients; reflect local inflammatory signaling from perivascular niches.
CXCL17, GDF-15, PTX3 (pentraxin-3) – recent proteomic hits; PTX3 more specific for vascular/adventitial inflammation in some heritable aortopathies.
TNF-α, IL-10, sST2 – useful for acute AD diagnosis/prognosis.
Acute injury / structural damage markers
D-dimer – high sensitivity for acute AD (cut-offs ~500 ng/mL); also tracks thrombus and growth.
Smooth-muscle myosin heavy chain (smMHC), calponin, soluble elastin fragments (sELAF) – early markers of medial/adventitial disruption.
Homocysteine, kynurenine, N-acetylneuraminic acid – metabolomic candidates from tissue-informed studies of acute AD.
Emerging proteomic / multi-omic panels (2020–2026)
CST3 (cystatin 3) + MMP12 + MEGF9 + CXCL17 – four-protein panel predicting future AA/AD risk (AUC improvement when combined with demographics; elevations detectable up to 10 years prior).
GDF-15, IL-6, CD58, LY9, Siglec-7 – validated for type-B AD.
NPC2, IGFBP7, ANG2 + Aggrecan panel, Lp(a), TMAO – associated with presence, growth rate, or surgical risk.
Additional circulating EV or plasma proteins: WFDC2, ACTN1, and various complement/coagulation factors.
Notes on external/perivascular tissue relevance
PVAT dysfunction and adventitial inflammation drive local release of IL-6, MCP-1, MMPs, and adipokines that appear in serum and correlate with wall remodeling and external support loss.
In FSI/growth-and-remodeling models, these markers can inform the time-dependent degradation of the ECM integrity variable (E) and the mechanical contribution of external tissue support.
Most markers remain investigational; none is yet standard for routine clinical decision-making. Calibration against longitudinal imaging and histology is required for quantitative coupling to continuum models.
Sources span narrative reviews, proteomic discovery studies, and cohort analyses published through mid-2026. Values and performance metrics (sensitivity/specificity/AUC) vary by cohort, assay, and endpoint; consult the original papers for exact cut-offs and multivariable adjustments.
